\section{Conclusions and Future Work}
\label{sec:conclusion}

We show that gradient-based optimization is viable for satellite constellation configuration when coverage and revisit metrics are made differentiable through continuous relaxations.
The relaxed objectives track their hard counterparts (\secref{sec:tightness}) and produce well-conditioned gradients where finite-difference approaches do not~\citep{paek_optimization_2019}.
On the Walker recovery testbed (\secref{sec:uniform}), the pipeline recovers the canonical Walker 24/6/1 hard metrics (\SI{40.3}{\percent} coverage, \SI{48}{\minute} mean worst-case revisit) within ${\sim}750$ iterations, whereas tuned simulated annealing, genetic algorithm, and differential evolution baselines plateau at \SI{60}{\minute}--\SI{74}{\minute} revisit with ${\sim}4\times$ the evaluation budget (\secref{sec:baselines}, \tabref{tab:baselines}).
On the regional Europe problem (\secref{sec:weighted}), the same pipeline escapes the circular-LEO basin and discovers Molniya-class highly elliptical orbits at the critical inclination---obtaining \SI{99.3}{\percent} coverage and \SI{3.7}{\minute} revisit versus \SI{10.0}{\percent} and \SI{247.5}{\minute} for a Walker 4/2/1 baseline---without any domain-specific inductive biases.

\paragraph{Scope and additional constraints.}
The scope of this work is the configuration problem: given a fixed satellite count $N$ and plane structure, the optimizer assigns orbital elements. Sizing (choosing $N$ and plane count) remains discrete and is handled outside the gradient loop. The free-geometry Exp 3 result illustrates that the unconstrained optimum can move far from the regimes most missions operate in: the discovered Molniya-class orbits traverse the inner Van Allen belt near \SI{10000}{\kilo\metre} apogee, with qualitatively different radiation, shielding, and duty-cycle implications than the LEO perigee. Rather than treating this as a limitation, we note that the same sigmoid reparameterization used for altitude in \secref{para:interval_constraints} admits any box-valued constraint on the orbital elements; the question is which constraints the mission imposes.

Physical-envelope constraints enter as interval bounds on the relevant element: bounding apogee below the inner-belt threshold (or avoiding the South Atlantic Anomaly for polar missions) becomes an interval constraint on $r_a = a(1+e)$ composable with the perigee reparameterization of \eqref{eq:reparam_ae}. Launch-feasibility constraints are similar in form: commercial rideshares deliver to a small set of inclinations, so matching a launch-available band is an interval constraint on $i$. A softer version is a differentiable $\Delta v$ penalty from a nominal dropoff orbit (e.g.\ \SI{500}{\kilo\metre} circular at the launch inclination) to the optimized elements, which constrains the optimizer to geometries the spacecraft can actually reach.

Mission-level visibility constraints enter through additional factors in the noisy-OR. Many imaging missions require Sun-elevation above the target, a second discontinuous condition parameterized by the solar ephemeris~\citep{vallado_fundamentals_2022}; the same soft-sigmoid treatment as satellite elevation composes into the noisy-OR as an extra factor, and the daylight-only coverage fraction is a product of two soft visibilities. Multi-pass requirements for stereo, change detection, or integration-time budgets need $K$ independent views per day rather than one; this is expressible as a count over soft visibilities followed by either a second LogSumExp or a $p$-norm soft-max when numerical conditioning becomes tight.

Longer propagation horizons would expose effects the \SI{24}{\hour} snapshot does not capture: $J_2$ RAAN drift, atmospheric drag, seasonal repeat tracks, and the polar-motion offset neglected by the simplified TEME-to-ECEF rotation of \eqref{eq:teme_to_ecef}. Beyond ${\sim}$\SI{7}{days} the accumulated polar-motion offset reaches the arcsecond scale (kilometer-level ground error), so the full IAU~2006/2000A precession--nutation chain would replace the single-axis rotation. Extending the pipeline in this direction is feasible either by parallel multi-epoch aggregation or by incorporating $J_2$ secular rates directly into the propagator, potentially letting the optimizer exploit RAAN drift or Sun-synchronous repeat tracks rather than treat them as nuisance physics.

The pipeline is otherwise agnostic to the objective: heterogeneous sensor placement, constellation reconfiguration, and multi-objective coverage with mission-specific elevation constraints are all expressible as weighted sums over the same graph.
